% GENERATED FILE. Edit canonical lessons, then run npm run generate:books.
% canonical-source-hash: fnv1a64:5bdd52dadd2bdb38
% canonical-lessons: LA-C291-primo, LA-C291-tunc, LA-C291-antequam, LA-C291-nondum, LA-C291-demum

\chapter{Prīmō, tunc, dēmum}
\label{ch:la-primo-tunc-demum}
\hlchaptermodality{291}

\begin{chapteropening}
\textbf{By the end of this chapter:} \emph{I can set the stages of a story in time: at first, at that time, before, not yet and at last.}

Complete the last lesson of chapter 291: I can set the stages of a story in time: at first, at that time, before, not yet and at last.
\end{chapteropening}

\section[prīmō]{prīmō — at first, at the beginning}

\label{lesson:LA-C291-primo}



\begin{quote}
Before the new one: say the Latin for a ditch, then the Latin for a rampart.
\end{quote}

\subsection*{You'll want to know: prīmō}
\textbf{prīmō} — \textquotedblleft{}at first, at the beginning\textquotedblright{}.

\begin{quote}
Prīmō caelum serēnum erat; deinde pluit.
\end{quote}

\begin{quote}
\emph{At first the sky was clear; then it rained.}
\end{quote}

An adverb. \emph{Prīmō} marks a first stage that later changes; \emph{prīmum} (taught) counts the steps of a list.

\begin{cousinweb}
\textbf{prīmō} is the ablative of \emph{prīmus}, \textquotedblleft{}first\textquotedblright{}, used as an adverb (Lewis \& Short): at the first point in time. English \emph{prime} and \emph{primary} come from \emph{prīmus}.
\end{cousinweb}

\begin{grammarlens}[title={What you've built}]
One more word for telling, explaining and describing at B1.
\end{grammarlens}

\subsection*{Your turn}
\begin{itemize}
\raggedright
  \item \emph{Say it:} \emph{prīmō}
  \item \emph{Say it:} \emph{prīmō}, once more
  \item \emph{Say it:} a sentence of your own with \emph{prīmō}
  \item \emph{Recall:} say the Latin for a ditch, then the Latin for a rampart, then say \emph{prīmō} again
\end{itemize}

\begin{tcolorbox}[breakable,colback=teal!4,colframe=teal!35!black,title={Before you move on}]
What does prīmō mean? (\textquotedblleft{}At first, at the beginning\textquotedblright{}.) Say it once more.
\end{tcolorbox}
\section[tunc]{tunc — then, at that time}

\label{lesson:LA-C291-tunc}



\begin{quote}
Before the new one: say the Latin for a rampart, then the Latin for at first.
\end{quote}

\subsection*{You'll want to know: tunc}
\textbf{tunc} — \textquotedblleft{}then, at that time\textquotedblright{}.

\begin{quote}
Tunc urbs parva erat, sed mūrī altī erant.
\end{quote}

\begin{quote}
\emph{At that time the city was small, but the walls were high.}
\end{quote}

A stronger \emph{tum} (taught): at that very time. Livy and later writers use it freely in stories.

\begin{cousinweb}
\textbf{tunc} is built on \emph{tum}, \textquotedblleft{}then\textquotedblright{}, with the pointing particle \emph{-ce} that is also seen in \emph{hic} (Lewis \& Short): then, at that point.
\end{cousinweb}

\begin{grammarlens}[title={What you've built}]
One more word for telling, explaining and describing at B1.
\end{grammarlens}

\subsection*{Your turn}
\begin{itemize}
\raggedright
  \item \emph{Say it:} \emph{tunc}
  \item \emph{Say it:} \emph{tunc}, once more
  \item \emph{Say it:} a sentence of your own with \emph{tunc}
  \item \emph{Recall:} say the Latin for a rampart, then the Latin for at first, then say \emph{tunc} again
\end{itemize}

\begin{tcolorbox}[breakable,colback=teal!4,colframe=teal!35!black,title={Before you move on}]
What does tunc mean? (\textquotedblleft{}Then, at that time\textquotedblright{}.) Say it once more.
\end{tcolorbox}
\section[antequam]{antequam — before (joining two clauses)}

\label{lesson:LA-C291-antequam}



\begin{quote}
Before the new one: say the Latin for at first, then the Latin for then.
\end{quote}

\subsection*{You'll want to know: antequam}
\textbf{antequam} — \textquotedblleft{}before (joining two clauses)\textquotedblright{}.

\begin{quote}
Antequam lūx vēnit, mīlitēs ē castrīs discessērunt.
\end{quote}

\begin{quote}
\emph{Before daylight came, the soldiers left the camp.}
\end{quote}

A conjunction, often split as \emph{ante … quam}. With the perfect for a plain past event.

\begin{cousinweb}
\textbf{antequam} joins \emph{ante}, \textquotedblleft{}before\textquotedblright{}, and \emph{quam}, \textquotedblleft{}than, as\textquotedblright{} (Lewis \& Short): earlier than the moment when.
\end{cousinweb}

\begin{grammarlens}[title={What you've built}]
One more word for telling, explaining and describing at B1.
\end{grammarlens}

\subsection*{Your turn}
\begin{itemize}
\raggedright
  \item \emph{Say it:} \emph{antequam}
  \item \emph{Say it:} \emph{antequam}, once more
  \item \emph{Say it:} a sentence of your own with \emph{antequam}
  \item \emph{Recall:} say the Latin for at first, then the Latin for then, then say \emph{antequam} again
\end{itemize}

\begin{tcolorbox}[breakable,colback=teal!4,colframe=teal!35!black,title={Before you move on}]
What does antequam mean? (\textquotedblleft{}Before (joining two clauses)\textquotedblright{}.) Say it once more.
\end{tcolorbox}
\section[nōndum]{nōndum — not yet}

\label{lesson:LA-C291-nondum}



\begin{quote}
Before the new one: say the Latin for then, then the Latin for before.
\end{quote}

\subsection*{You'll want to know: nōndum}
\textbf{nōndum} — \textquotedblleft{}not yet\textquotedblright{}.

\begin{quote}
Nūntius nōndum vēnerat; itaque exspectābāmus.
\end{quote}

\begin{quote}
\emph{The messenger had not yet come, so we were waiting.}
\end{quote}

The opposite of \emph{iam} (taught), 'already'. \emph{Nōndum} + pluperfect sets a scene in which something is still awaited.

\begin{cousinweb}
\textbf{nōndum} joins \emph{nōn}, \textquotedblleft{}not\textquotedblright{}, and \emph{dum}, \textquotedblleft{}yet, still\textquotedblright{} (Lewis \& Short): not as yet.
\end{cousinweb}

\begin{grammarlens}[title={What you've built}]
One more word for telling, explaining and describing at B1.
\end{grammarlens}

\subsection*{Your turn}
\begin{itemize}
\raggedright
  \item \emph{Say it:} \emph{nōndum}
  \item \emph{Say it:} \emph{nōndum}, once more
  \item \emph{Say it:} a sentence of your own with \emph{nōndum}
  \item \emph{Recall:} say the Latin for then, then the Latin for before, then say \emph{nōndum} again
\end{itemize}

\begin{tcolorbox}[breakable,colback=teal!4,colframe=teal!35!black,title={Before you move on}]
What does nōndum mean? (\textquotedblleft{}Not yet\textquotedblright{}.) Say it once more.
\end{tcolorbox}
\section[dēmum]{dēmum — at last, only then}

\label{lesson:LA-C291-demum}



\begin{quote}
Before the new one: say the Latin for before, then the Latin for not yet.
\end{quote}

\subsection*{You'll want to know: dēmum}
\textbf{dēmum} — \textquotedblleft{}at last, only then\textquotedblright{}.

\begin{quote}
Post multās hōrās dēmum ad urbem pervēnimus.
\end{quote}

\begin{quote}
\emph{After many hours we at last reached the city.}
\end{quote}

\emph{Tum dēmum}: only then. It stresses that something came late, after waiting; close to \emph{tandem} (taught).

\begin{cousinweb}
\textbf{dēmum} is explained by Lewis \& Short as a superlative-like form built on \emph{dē}, \textquotedblleft{}down from\textquotedblright{}: the lowest, last point. The details are not settled.
\end{cousinweb}

\begin{grammarlens}[title={What you've built}]
One more word for telling, explaining and describing at B1.
\end{grammarlens}

\subsection*{Your turn}
\begin{itemize}
\raggedright
  \item \emph{Say it:} \emph{dēmum}
  \item \emph{Say it:} \emph{dēmum}, once more
  \item \emph{Say it:} a sentence of your own with \emph{dēmum}
  \item \emph{Recall:} say the Latin for before, then the Latin for not yet, then say \emph{dēmum} again
\end{itemize}

\begin{tcolorbox}[breakable,colback=teal!4,colframe=teal!35!black,title={Before you move on}]
What does dēmum mean? (\textquotedblleft{}At last, only then\textquotedblright{}.) Say it once more.
\end{tcolorbox}
